\documentclass[10pt,a4paper]{amsart}
\usepackage{amsmath,amssymb,amsthm,graphicx}
\usepackage{verbatim}
\newenvironment{reference}{\comment}{\endcomment}
%\newenvironment{reference}{}{}
\newenvironment{slogan}{\comment}{\endcomment}
\newenvironment{history}{\comment}{\endcomment}
% For commutative diagrams we use Xy-pic
\usepackage[all]{xy}

% We use 2cell for 2-commutative diagrams.
\xyoption{2cell}
% Theorem environments.
%
\theoremstyle{plain}
\newtheorem{theorem}[subsection]{Theorem}
\newtheorem{proposition}[subsection]{Proposition}
\newtheorem{lemma}[subsection]{Lemma}

\theoremstyle{definition}
\newtheorem{definition}[subsection]{Definition}
\newtheorem{example}[subsection]{Example}
\newtheorem{exercise}[subsection]{Exercise}
\newtheorem{situation}[subsection]{Situation}

\theoremstyle{remark}
\newtheorem{remark}[subsection]{Remark}
\newtheorem{remarks}[subsection]{Remarks}

\numberwithin{equation}{subsection}

% Macros
%
\def\lim{\mathop{\mathrm{lim}}\nolimits}
\def\colim{\mathop{\mathrm{colim}}\nolimits}
\def\Spec{\mathop{\mathrm{Spec}}}
\def\Hom{\mathop{\mathrm{Hom}}\nolimits}
\def\Ext{\mathop{\mathrm{Ext}}\nolimits}
\def\SheafHom{\mathop{\mathcal{H}\!\mathit{om}}\nolimits}
\def\SheafExt{\mathop{\mathcal{E}\!\mathit{xt}}\nolimits}
\def\Sch{\mathit{Sch}}
\def\Mor{\mathop{\mathrm{Mor}}\nolimits}
\def\Ob{\mathop{\mathrm{Ob}}\nolimits}
\def\Sh{\mathop{\mathit{Sh}}\nolimits}
\def\NL{\mathop{N\!L}\nolimits}
\def\CH{\mathop{\mathrm{CH}}\nolimits}
\def\proetale{{pro\text{-}\acute{e}tale}}
\def\etale{{\acute{e}tale}}
\def\QCoh{\mathit{QCoh}}
\def\Ker{\mathop{\mathrm{Ker}}}
\def\Im{\mathop{\mathrm{Im}}}
\def\Coker{\mathop{\mathrm{Coker}}}
\def\Coim{\mathop{\mathrm{Coim}}}

% Boxtimes
%
\DeclareMathSymbol{\boxtimes}{\mathbin}{AMSa}{"02}

%
% Macros for moduli stacks/spaces
%
\def\QCohstack{\mathcal{QC}\!\mathit{oh}}
\def\Cohstack{\mathcal{C}\!\mathit{oh}}
\def\Spacesstack{\mathcal{S}\!\mathit{paces}}
\def\Quotfunctor{\mathrm{Quot}}
\def\Hilbfunctor{\mathrm{Hilb}}
\def\Curvesstack{\mathcal{C}\!\mathit{urves}}
\def\Polarizedstack{\mathcal{P}\!\mathit{olarized}}
\def\Complexesstack{\mathcal{C}\!\mathit{omplexes}}
% \Pic is the operator that assigns to X its picard group, usage \Pic(X)
% \Picardstack_{X/B} denotes the Picard stack of X over B
% \Picardfunctor_{X/B} denotes the Picard functor of X over B
\def\Pic{\mathop{\mathrm{Pic}}\nolimits}
\def\Picardstack{\mathcal{P}\!\mathit{ic}}
\def\Picardfunctor{\mathrm{Pic}}
\def\Deformationcategory{\mathcal{D}\!\mathit{ef}}
\usepackage{cleveref}
\crefname{originallemma}{original lemma}{original lemmas}
\Crefname{originallemma}{Original lemma}{Original lemmas}
\makeatletter
\newlabel{lemma-cap-projective-bundle}{{02TW}{source}{}{sourcepointer-1}{}}
\newlabel{lemma-cap-projective-bundle@cref}{{[originallemma][0][]02TW}{[1][0][]source}{}{}{}}
\makeatother
\crefname{originallemma}{original lemma}{original lemmas}
\Crefname{originallemma}{Original lemma}{Original lemmas}
\makeatletter
\newlabel{lemma-flat-pullback-proper-pushforward}{{02RG}{source}{}{sourcepointer-2}{}}
\newlabel{lemma-flat-pullback-proper-pushforward@cref}{{[originallemma][0][]02RG}{[1][0][]source}{}{}{}}
\makeatother
\crefname{originallemma}{original lemma}{original lemmas}
\Crefname{originallemma}{Original lemma}{Original lemmas}
\makeatletter
\newlabel{lemma-pullback-affine-fibres-surjective}{{02TT}{source}{}{sourcepointer-3}{}}
\newlabel{lemma-pullback-affine-fibres-surjective@cref}{{[originallemma][0][]02TT}{[1][0][]source}{}{}{}}
\makeatother
\crefname{originallemma}{original lemma}{original lemmas}
\Crefname{originallemma}{Original lemma}{Original lemmas}
\makeatletter
\newlabel{lemma-pushforward-cap-c1}{{02SU}{source}{}{sourcepointer-4}{}}
\newlabel{lemma-pushforward-cap-c1@cref}{{[originallemma][0][]02SU}{[1][0][]source}{}{}{}}
\makeatother
\crefname{originallemma}{original lemma}{original lemmas}
\Crefname{originallemma}{Original lemma}{Original lemmas}
\makeatletter
\newlabel{lemma-relative-effective-cartier}{{02TR}{source}{}{sourcepointer-5}{}}
\newlabel{lemma-relative-effective-cartier@cref}{{[originallemma][0][]02TR}{[1][0][]source}{}{}{}}
\makeatother
\crefname{originallemma}{original lemma}{original lemmas}
\Crefname{originallemma}{Original lemma}{Original lemmas}
\makeatletter
\newlabel{lemma-restrict-to-open}{{02RX}{source}{}{sourcepointer-6}{}}
\newlabel{lemma-restrict-to-open@cref}{{[originallemma][0][]02RX}{[1][0][]source}{}{}{}}
\makeatother
\crefname{originallemma}{original lemma}{original lemmas}
\Crefname{originallemma}{Original lemma}{Original lemmas}
\makeatletter
\newlabel{situation-setup}{{02QL}{source}{}{sourcepointer-7}{}}
\newlabel{situation-setup@cref}{{[originallemma][0][]02QL}{[1][0][]source}{}{}{}}
\makeatother
\title{The projective space bundle formula: the Chow group of a projective bundle as a direct sum over the base}
\author{The Stacks Project Authors}
\date{stacks-project:chow.tex:local-snapshot-20261012; GNU FDL 1.2 or later}
\begin{document}
\maketitle
\noindent\textit{Editorial note.} All delivered mathematical text is the author's own, copied from the pinned Stacks Project chapter file chow.tex; no statement or proof step was rewritten, reordered or removed.\par\medskip
\noindent\textit{Editorial note.} Preamble provenance: the pinned chapter file opens with \textbackslash{}input\{preamble\} and carries no preamble text of its own, so the project's single Stacks-wide preamble staged at sources/stacks-chow/preamble.tex is excerpted verbatim as three spans, lines 8-12 (the verbatim package together with the project's own reference, slogan and history display-environment definitions, including the author's commented-out older form of the reference environment), lines 14-18 (Xy-pic setup) and lines 154-217 (the native theorem declarations and mathematical macros). The lines 8-12 span is required because a counted statement or its surrounding text may use those display environments, which the project typesets as comments. Nothing was added, removed or redefined; the wrapper class and the amsmath/amssymb/amsthm/graphicx base packages are supplied by the assembler.\par\medskip
\noindent\textit{Editorial note.} Same-chapter pointers are located at their \textbackslash{}label lines inside the pinned chapter file; cross-chapter pointers are located in the project's own tag list staged in the snapshot (sources/stacks-chow/tags.txt), at the verbatim lines carrying their canonical tags.\par\medskip
\noindent\textit{Allowed standard prerequisites (editorial).}\begin{itemize}
\item The pushforward of a power of the first Chern class of the twisting sheaf. Same chapter chow.tex, label lemma-cap-projective-bundle, invoked at proof line 6960. Gives that pushing forward c\_1(O\_P(1))\textasciicircum{}s cap pi\textasciicircum{}*alpha along the bundle is zero for s below r - 1 and returns alpha for s equal to r - 1, which is exactly the injectivity computation. The lemma is claimed separately as S14-algebra-cap-projective-bundle.
\item Flat pullback versus proper pushforward and the projection formula for capping with a first Chern class. Same chapter chow.tex, labels lemma-flat-pullback-proper-pushforward (proof line 6992) and lemma-pushforward-cap-c1 (proof line 6994). Let the surjectivity statement be checked component by component over the disjoint union of the irreducible components of the base, and lets the capping operation be moved through the flat base change.
\item Restriction to an open and the surjectivity of pullback along a morphism with affine space fibres. Same chapter chow.tex, labels lemma-restrict-to-open (proof line 7016), lemma-pullback-affine-fibres-surjective (proof lines 7043 and 7048) and lemma-relative-effective-cartier (proof line 7052). The first two let the base be replaced by a nonempty open and supply the successive lifts from the strata of the projective space, and the third identifies the class of the closure of a stratum with the cap of the first Chern class against the pulled back class.
\end{itemize}\medskip
\ifdefined\hypertarget\hypertarget{sourcepointer-1}{}\fi
\noindent\textit{Source reference.} Original lemma 02TW: sources/stacks-chow/pinned-original-source.tex, source lines 6868--6868. Navigation pointer for \textbackslash{}ref\{lemma-cap-projective-bundle\} used inside the counted span (line 6960); the argument that consumes it is retained in the span.\par\medskip
\ifdefined\hypertarget\hypertarget{sourcepointer-2}{}\fi
\noindent\textit{Source reference.} Original lemma 02RG: sources/stacks-chow/pinned-original-source.tex, source lines 2401--2401. Navigation pointer for \textbackslash{}ref\{lemma-flat-pullback-proper-pushforward\} used inside the counted span (line 6992); the argument that consumes it is retained in the span.\par\medskip
\ifdefined\hypertarget\hypertarget{sourcepointer-3}{}\fi
\noindent\textit{Source reference.} Original lemma 02TT: sources/stacks-chow/pinned-original-source.tex, source lines 5651--5651. Navigation pointer for \textbackslash{}ref\{lemma-pullback-affine-fibres-surjective\} used inside the counted span (line 7043); the argument that consumes it is retained in the span.\par\medskip
\ifdefined\hypertarget\hypertarget{sourcepointer-4}{}\fi
\noindent\textit{Source reference.} Original lemma 02SU: sources/stacks-chow/pinned-original-source.tex, source lines 4631--4631. Navigation pointer for \textbackslash{}ref\{lemma-pushforward-cap-c1\} used inside the counted span (line 6994); the argument that consumes it is retained in the span.\par\medskip
\ifdefined\hypertarget\hypertarget{sourcepointer-5}{}\fi
\noindent\textit{Source reference.} Original lemma 02TR: sources/stacks-chow/pinned-original-source.tex, source lines 5598--5598. Navigation pointer for \textbackslash{}ref\{lemma-relative-effective-cartier\} used inside the counted span (line 7052); the argument that consumes it is retained in the span.\par\medskip
\ifdefined\hypertarget\hypertarget{sourcepointer-6}{}\fi
\noindent\textit{Source reference.} Original lemma 02RX: sources/stacks-chow/pinned-original-source.tex, source lines 2949--2949. Navigation pointer for \textbackslash{}ref\{lemma-restrict-to-open\} used inside the counted span (line 7016); the argument that consumes it is retained in the span.\par\medskip
\ifdefined\hypertarget\hypertarget{sourcepointer-7}{}\fi
\noindent\textit{Source reference.} Original lemma 02QL: sources/stacks-chow/pinned-original-source.tex, source lines 1372--1372. Navigation pointer for \textbackslash{}ref\{situation-setup\} used inside the counted span (line 6933); the argument that consumes it is retained in the span.\par\medskip
\begin{lemma}[Projective space bundle formula]
\label{lemma-chow-ring-projective-bundle}
Let $(S, \delta)$ be as in Situation \ref{situation-setup}.
Let $X$ be locally of finite type over $S$.
Let $\mathcal{E}$ be a finite locally free $\mathcal{O}_X$-module
$\mathcal{E}$ of rank $r$. Let $(\pi : P \to X, \mathcal{O}_P(1))$
be the projective bundle associated to $\mathcal{E}$.
The map
$$
\bigoplus\nolimits_{i = 0}^{r - 1}
\CH_{k + i}(X)
\longrightarrow
\CH_{k + r - 1}(P),
$$
$$
(\alpha_0, \ldots, \alpha_{r-1})
\longmapsto
\pi^*\alpha_0 +
c_1(\mathcal{O}_P(1)) \cap \pi^*\alpha_1
+ \ldots +
c_1(\mathcal{O}_P(1))^{r - 1} \cap \pi^*\alpha_{r-1}
$$
is an isomorphism.
\end{lemma}

\begin{proof}
Fix $k \in \mathbf{Z}$. We first show the map is injective.
Suppose that $(\alpha_0, \ldots, \alpha_{r - 1})$ is an element
of the left hand side that maps to zero.
By Lemma \ref{lemma-cap-projective-bundle} we see that
$$
0 = \pi_*(\pi^*\alpha_0 +
c_1(\mathcal{O}_P(1)) \cap \pi^*\alpha_1
+ \ldots +
c_1(\mathcal{O}_P(1))^{r - 1} \cap \pi^*\alpha_{r-1})
= \alpha_{r - 1}
$$
Next, we see that
$$
0 = \pi_*(c_1(\mathcal{O}_P(1)) \cap (\pi^*\alpha_0 +
c_1(\mathcal{O}_P(1)) \cap \pi^*\alpha_1
+ \ldots +
c_1(\mathcal{O}_P(1))^{r - 2} \cap \pi^*\alpha_{r - 2}))
= \alpha_{r - 2}
$$
and so on. Hence the map is injective.

\medskip\noindent
It remains to show the map is surjective.
Let $X_i$, $i \in I$ be the irreducible components of $X$.
Then $P_i = \mathbf{P}(\mathcal{E}|_{X_i})$, $i \in I$
are the irreducible components of $P$. Consider the commutative
diagram
$$
\xymatrix{
\coprod P_i \ar[d]_{\coprod \pi_i} \ar[r]_p & P \ar[d]^\pi \\
\coprod X_i \ar[r]^q & X
}
$$
Observe that $p_*$ is surjective. If $\beta \in \CH_k(\coprod X_i)$
then $\pi^* q_* \beta = p_*(\coprod \pi_i)^* \beta$, see
Lemma \ref{lemma-flat-pullback-proper-pushforward}. Similarly for
capping with $c_1(\mathcal{O}(1))$ by
Lemma \ref{lemma-pushforward-cap-c1}.
Hence, if the map of the lemma is surjective for each
of the morphisms $\pi_i : P_i \to X_i$, then the map is
surjective for $\pi : P \to X$. Hence we may assume $X$ is irreducible.
Thus $\dim_\delta(X) < \infty$ and in particular we may use
induction on $\dim_\delta(X)$.

\medskip\noindent
The result is clear if $\dim_\delta(X) < k$.
Let $\alpha \in \CH_{k + r - 1}(P)$.
For any locally closed subscheme $T \subset X$ denote
$\gamma_T : \bigoplus \CH_{k + i}(T) \to \CH_{k + r - 1}(\pi^{-1}(T))$
the map
$$
\gamma_T(\alpha_0, \ldots, \alpha_{r - 1})
= \pi^*\alpha_0 + \ldots +
c_1(\mathcal{O}_{\pi^{-1}(T)}(1))^{r - 1} \cap \pi^*\alpha_{r - 1}.
$$
Suppose for some nonempty open $U \subset X$ we have
$\alpha|_{\pi^{-1}(U)} = \gamma_U(\alpha_0, \ldots, \alpha_{r - 1})$.
Then we may choose lifts $\alpha'_i \in \CH_{k + i}(X)$ and we
see that $\alpha - \gamma_X(\alpha'_0, \ldots, \alpha'_{r - 1})$
is by Lemma \ref{lemma-restrict-to-open}
rationally equivalent to a $k$-cycle on $P_Y = \mathbf{P}(\mathcal{E}|_Y)$
where $Y = X \setminus U$ as a reduced closed subscheme.
Note that $\dim_\delta(Y) < \dim_\delta(X)$.
By induction the result holds
for $P_Y \to Y$ and hence the result holds for $\alpha$.
Hence we may replace $X$ by any nonempty open of $X$.

\medskip\noindent
In particular we may assume that $\mathcal{E} \cong \mathcal{O}_X^{\oplus r}$.
In this case $\mathbf{P}(\mathcal{E}) = X \times \mathbf{P}^{r - 1}$.
Let us use the stratification
$$
\mathbf{P}^{r - 1} = \mathbf{A}^{r - 1}
\amalg \mathbf{A}^{r - 2}
\amalg \ldots
\amalg \mathbf{A}^0
$$
The closure of each stratum is a $\mathbf{P}^{r - 1 - i}$ which is a
representative of $c_1(\mathcal{O}(1))^i \cap [\mathbf{P}^{r - 1}]$.
Hence $P$ has a similar stratification
$$
P = U^{r - 1} \amalg U^{r - 2} \amalg \ldots \amalg U^0
$$
Let $P^i$ be the closure of $U^i$. Let $\pi^i : P^i \to X$
be the restriction of $\pi$ to $P^i$.
Let $\alpha \in \CH_{k + r - 1}(P)$. By
Lemma \ref{lemma-pullback-affine-fibres-surjective}
we can write $\alpha|_{U^{r - 1}} = \pi^*\alpha_0|_{U^{r - 1}}$
for some $\alpha_0 \in \CH_k(X)$. Hence the difference
$\alpha - \pi^*\alpha_0$ is the image of some
$\alpha' \in \CH_{k + r - 1}(P^{r - 2})$.
By Lemma \ref{lemma-pullback-affine-fibres-surjective}
again we can write
$\alpha'|_{U^{r - 2}} = (\pi^{r - 2})^*\alpha_1|_{U^{r - 2}}$
for some $\alpha_1 \in \CH_{k + 1}(X)$.
By Lemma \ref{lemma-relative-effective-cartier}
we see that the image of $(\pi^{r - 2})^*\alpha_1$
represents $c_1(\mathcal{O}_P(1)) \cap \pi^*\alpha_1$.
We also see that
$\alpha - \pi^*\alpha_0 - c_1(\mathcal{O}_P(1)) \cap \pi^*\alpha_1$
is the image of some $\alpha'' \in \CH_{k + r - 1}(P^{r - 3})$.
And so on.
\end{proof}

\end{document}
